\chapter{实例方案——《函数的性质——函数的凹凸性》习题课}

函数具有高中数学的内容基础性，有着不可替代的重要位置。函数的性质是反应函数变化规律的重要研究途径，学生应该掌握函数的一般性质，并且学会自主研究函数的其他性质，因此该问题具有探究性。本案例采用课本中习题进行探究，习题选自新人教A版高中数学必修一P101第8题。在教学中通过对课本习题进行延伸、扩展，在讲解习题时适当对习题进行思维拓展，积极引导学生对原习题进行变式，经过问题假设、归纳证明等再发现的思维活动，探索变式的可行性与一般化方法，促进了学生的探索以及创新能力。

课前要求：请学生利用网络资源或书本资源搜集有关凹凸函数的资料，可以包括定义、图象、性质等等。

教学过程：

证明：

\begin{enumerate}[(1)]
  \item 存 $f(x) = a x + b$, 则 $f\left(\frac{x_{1}+x_{2}}{2}\right)=\frac{f\left(x_{1}\right)+f\left(x_{2}\right)}{2}$;
  \item 若 $g(x)=x^{2}+a x+b$, 则 $g\left(\frac{x_{1}+x_{2}}{2}\right) \leq \frac{g\left(x_{1}\right)+g\left(x_{2}\right)}{2}$.
\end{enumerate}

师：这一章我们学习了函数的概念以及性质，函数是反应事物变化规律的重要模型，其性质是我们用于研究函数的重要途径。本道习题考查的也是函数性质的应用，同学们容易发现这道习题与我们所学习的函数性质不大相同，面对这个习题我们应该如何下手呢？它与我们学习过的哪一个性质相似呢？你能想到哪些性质与这道习题相关？

生：结合教师提问的方向进行同桌讨论与思考，容易发现这与函数的单调性在结构上有些相似，该问题与函数的求值相关。

（教师请同学们动笔操作，并且请学生进行板演，教师详细讲解证明过程，并针对本习题提出拓展性的探究问题。）

师：根据（2）的结论，你能否改变（2）中的条件，探究其结论是否正确？或者改变结论并证明。你能有几种不同的条件？能否将该习题进行结论推广？

（教师引导学生进行习题推广的探究活动，通过改变条件证明结论或证明不同的结论，大胆放手让学生尝试。安排学生前后桌四人为一小组，合作探究，并让学生板演探究结果。）

生：可能会有以下几种拓展：

\begin{enumerate}[a)]
  \item 若 $g(x)=2 x^{2}+a x+b$, 则 $g\left(\frac{x_{1}+x_{2}}{2}\right) \leq \frac{g\left(x_{1}\right)+g\left(x_{2}\right)}{2}$
  \item 若 $g(x)=-x^{2}+a x+b$, 则 $g\left(\frac{x_{1}+x_{2}}{2}\right) \leq \frac{g\left(x_{1}\right)+g\left(x_{2}\right)}{2}$
  \item 若 $g(x)=a x^{2}+b x+c(a \neq 0)$, 则 $g\left(\frac{x_{1}+x_{2}}{2}\right) \leq \frac{g\left(x_{1}\right)+g\left(x_{2}\right)}{2}$
  \item 若 $g(x)=x^{2}+a x+b$, 则 $g\left(\frac{x_{1}+x_{2}+x_{3}}{3}\right) \leq \frac{g\left(x_{1}\right)+g\left(x_{2}\right)+g\left(x_{3}\right)}{3}$
  \item 若 $g(x)=x^{2}+ax+b$, 则 $g\left(\frac{x_{1}+x_{2}+x_{3}+\cdots x_{n}}{n}\right) \leq \frac{g\left(x_{1}\right)+g\left(x_{2}\right)+g\left(x_{3}\right)+\cdots g\left(x_{n}\right)}{n}$
  \item 若 $g(x)=a^{x}(a>1)$, 则 $g\left(\frac{x_{1}+x_{2}}{2}\right) \leq \frac{g\left(x_{1}\right)+g\left(x_{2}\right)}{2}$; 若 $0<a<1$, 则 $g\left(\frac{x_{1}+x_{2}}{2}\right) \geq \frac{g\left(x_{1}\right)+g\left(x_{2}\right)}{2}$
  \item 若 $g(x)=\log _{a} x(a>1)$, 则 $g\left(\frac{x_{1}+x_{2}}{2}\right) \geq \frac{g\left(x_{1}\right)+g\left(x_{2}\right)}{2}$; 若 $0<a<1$, 则 $g\left(\frac{x_{1}+x_{2}}{2}\right) \leq \frac{g\left(x_{1}\right)+g\left(x_{2}\right)}{2}$
\end{enumerate}

（教师应肯定学生的探究方向，学生发散思维对条件和结论进行改变推广，提出了多种多样的变式，接下来教师应当引导学生进行命题的证明。）

师：这些探究方向都非常好。$a$和$b$的想法类似，改变条件中的二次函数的系数；$c$将特殊二次函数推广到一般的二次函数，对二次项系数做了不同的处理；d则是对结论进行从2到3的推广，$e$甚至推广到$n$的情况，更加一般化；$f$和$g$将函数变为指对数函数进行探究……下面请同学们对这些命题进行严格的证明。

（拥有原命题的证明做铺垫，这些推论的证明对学生不难，重要的是发现和探索命题的过程。教师引导学生进行命题的拓展，保留正确的命题，修改错误的命题。）

师：根据同学们的证明，我们发现$a$、$d$、$e$是正确的，$b$、$c$、$f$、$g$错误，针对这些错误的命题，同学们能否进行修正？

生：基于$a$、$b$、$c$这三个命题的真伪，我们可以发现二次项系数的正负影响了结论的正误，应改为若$g(x)=a x^{2}+b x+c(a>0)$, 则 $g\left(\frac{x_{1}+x_{2}}{2}\right) \leq \frac{g(x_{1})+g(x_{2})}{2}$; 若 $a<0$, 则 $g\left(\frac{x_{1}+x_{2}}{2}\right) \geq \frac{g(x_{1})+g(x_{2})}{2}$ 。对于指对数函数，他们有着不同的性质，不能等同起来，$f$改为若$g(x)=a^{x}(a>0$ 且 $a \neq 1)$，都有$g\left(\frac{x_{1}+x_{2}}{2}\right) \leq \frac{g(x_{1})+g(x_{2})}{2}$；$g$缺少了自变量的取值范围，应该添加。

（通过对比指对数函数命题的异同点，结合前期学生自主查询凹凸函数的资料，教师可以适时抛出这些问题与凹凸函数的关联，积极引导学生总结规律，得出性质，并对今日所学作出总结和归纳，提出展望。教师板书一般步骤，课堂小结也可由学生总结回答。）

师：结合指对数函数的命题结论，你能发现什么规律吗？满足这些结论的函数叫什么名字？结合所查到的资料，你能总结出该类函数的性质吗？

生：凹凸函数的性质：如果$f(x)$是定义在区间 $D$ 上的下凸函数，则 $\forall x_{1}, x_{2} \in D$, 有 $f\left(\frac{x_{1}+x_{2}}{2}\right) \leq$ $\frac{f(x_{1})+f(x_{2})}{2}$; 如果 $f(x)$ 是定义在区间 $D$ 上的上凸函数, 则 $\forall x_{1}, x_{2} \in D$, 有 $f\left(\frac{x_{1}+x_{2}}{2}\right) \geq \frac{f(x_{1})+f(x_{2})}{2}$ 。

师：凹凸函数的定义与性质不止今天所学的这些，其定义和图象有待同学们深入学习，继续探索。今天我们经历了数学家的科学发现，所有命题、结论的发现都依赖于数学家们细致的观察与发现，进行猜想再验证解决。因此，在碰到同类问题时，甚至是完全陌生的习题，我们也可以采用探究性学习的一般步骤对习题进行研究，挖掘其更深刻的内涵。